\documentclass{anstrans}
%%%%%%%%%%%%%%%%%%%%%%%%%%%%%%%%%%%
\title{LEU+ to HALEU Transitions in Advanced Reactor Fuel Cycles}
\author{Nathan S. Ryan,$^{*}$ Kathryn D. Huff,$^{*}$ and Madicken Munk,$^{*, \dagger}$}

\institute{
$^{*}$Advanced Reactors and Fuel Cycles Group, University of Illinois,
Urbana, IL, nsryan2@illinois.edu
\and
$^{\dagger}$Scientific Computing, Reactor Analysis and Modeling Group, Oregon State University, Corvallis, OR
}

% Optional disclaimer: remove this command to hide
% \disclaimer{Notice: this manuscript is a work of fiction. Any resemblance to
% actual articles, living or dead, is purely coincidental.}

%%%% packages and definitions (optional)
\usepackage{graphicx} % allows inclusion of graphics
\usepackage{booktabs} % nice rules (thick lines) for tables
\usepackage{microtype} % improves typography for PDF
\usepackage[acronym,toc]{glossaries}
%\newacronym{<++>}{<++>}{<++>}
\newacronym[longplural={metric tons of heavy metal}]{MTHM}{MTHM}{metric ton of heavy metal}
\newacronym{ABM}{ABM}{agent-based modeling}
\newacronym{ACDIS}{ACDIS}{Program in Arms Control \& Domestic and International Security}
\newacronym{AHTR}{AHTR}{Advanced High Temperature Reactor}
\newacronym{ANDRA}{ANDRA}{Agence Nationale pour la gestion des D\'echets RAdioactifs, the French National Agency for Radioactive Waste Management}
\newacronym{ANL}{ANL}{Argonne National Laboratory}
\newacronym{API}{API}{application programming interface}
\newacronym{ARE}{ARE}{Aircraft Reactor Experiment}
\newacronym{ARFC}{ARFC}{Advanced Reactors and Fuel Cycles}
\newacronym{ASME}{ASME}{American Society of Mechanical Engineers}
\newacronym{ATWS}{ATWS}{Anticipated Transient Without Scram}
\newacronym{BDBE}{BDBE}{Beyond Design Basis Event}
\newacronym{BIDS}{BIDS}{Berkeley Institute for Data Science}
\newacronym{CAFCA}{CAFCA}{ Code for Advanced Fuel Cycles Assessment }
\newacronym{CDTN}{CDTN}{Centro de Desenvolvimento da Tecnologia Nuclear}
\newacronym{CEA}{CEA}{Commissariat \`a l'\'Energie Atomique et aux \'Energies Alternatives}
\newacronym{CI}{CI}{continuous integration}
\newacronym{CNEN}{CNEN}{Comiss\~{a}o Nacional de Energia Nuclear}
\newacronym{CNERG}{CNERG}{Computational Nuclear Engineering Research Group}
\newacronym{COSI}{COSI}{Commelini-Sicard}
\newacronym{COTS}{COTS}{commercial, off-the-shelf}
\newacronym{CSNF}{CSNF}{commercial spent nuclear fuel}
\newacronym{CTAH}{CTAHs}{Coiled Tube Air Heaters}
\newacronym{CUBIT}{CUBIT}{CUBIT Geometry and Mesh Generation Toolkit}
\newacronym{CURIE}{CURIE}{Centralized Used Fuel Resource for Information Exchange}
\newacronym{DAG}{DAG}{directed acyclic graph}
\newacronym{DANESS}{DANESS}{Dynamic Analysis of Nuclear Energy System Strategies}
\newacronym{DBE}{DBE}{Design Basis Event}
\newacronym{DESAE}{DESAE}{Dynamic Analysis of Nuclear Energy Systems Strategies}
\newacronym{DHS}{DHS}{Department of Homeland Security}
\newacronym{DOE}{DOE}{Department of Energy}
\newacronym{DRACS}{DRACS}{Direct Reactor Auxiliary Cooling System}
\newacronym{DRE}{DRE}{dynamic resource exchange}
\newacronym{DSNF}{DSNF}{DOE spent nuclear fuel}
\newacronym{DYMOND}{DYMOND}{Dynamic Model of Nuclear Development }
\newacronym{EBS}{EBS}{Engineered Barrier System}
\newacronym{EDZ}{EDZ}{Excavation Disturbed Zone}
\newacronym{EIA}{EIA}{U.S. Energy Information Administration}
\newacronym{EPA}{EPA}{Environmental Protection Agency}
\newacronym{EP}{EP}{Engineering Physics}
\newacronym{FCO}{FCO}{Fuel Cycle Options}
\newacronym{FCT}{FCT}{Fuel Cycle Technology}
\newacronym{FEHM}{FEHM}{Finite Element Heat and Mass Transfer}
\newacronym{FEPs}{FEPs}{Features, Events, and Processes}
\newacronym{FHR}{FHR}{Fluoride-Salt-Cooled High-Temperature Reactor}
\newacronym{FLiBe}{FLiBe}{Fluoride-Lithium-Beryllium}
\newacronym{GDSE}{GDSE}{Generic Disposal System Environment}
\newacronym{GDSM}{GDSM}{Generic Disposal System Model}
\newacronym{GENIUSv1}{GENIUSv1}{Global Evaluation of Nuclear Infrastructure Utilization Scenarios, Version 1}
\newacronym{GENIUSv2}{GENIUSv2}{Global Evaluation of Nuclear Infrastructure Utilization Scenarios, Version 2}
\newacronym{GENIUS}{GENIUS}{Global Evaluation of Nuclear Infrastructure Utilization Scenarios}
\newacronym{GPAM}{GPAM}{Generic Performance Assessment Model}
\newacronym{GRSAC}{GRSAC}{Graphite Reactor Severe Accident Code}
\newacronym{GUI}{GUI}{graphical user interface}
\newacronym{HLW}{HLW}{high level waste}
\newacronym{HPC}{HPC}{high-performance computing}
\newacronym{HTC}{HTC}{high-throughput computing}
\newacronym{HTGR}{HTGR}{High Temperature Gas-Cooled Reactor}
\newacronym{IAEA}{IAEA}{International Atomic Energy Agency}
\newacronym{IEMA}{IEMA}{Illinois Emergency Mangament Agency}
\newacronym{INL}{INL}{Idaho National Laboratory}
\newacronym{IPRR1}{IRP-R1}{Instituto de Pesquisas Radioativas Reator 1}
\newacronym{IRP}{IRP}{Integrated Research Project}
\newacronym{ISFSI}{ISFSI}{Independent Spent Fuel Storage Installation}
\newacronym{ISRG}{ISRG}{Independent Student Research Group}
\newacronym{JFNK}{JFNK}{Jacobian-Free Newton Krylov}
\newacronym{LANL}{LANL}{Los Alamos National Laboratory}
\newacronym{LBNL}{LBNL}{Lawrence Berkeley National Laboratory}
\newacronym{LCOE}{LCOE}{levelized cost of electricity}
\newacronym{LDRD}{LDRD}{laboratory directed research and development}
\newacronym{LFR}{LFR}{Lead-Cooled Fast Reactor}
\newacronym{LLNL}{LLNL}{Lawrence Livermore National Laboratory}
\newacronym{LMFBR}{LMFBR}{Liquid Metal Fast Breeder Reactor}
\newacronym{LOFC}{LOFC}{Loss of Forced Cooling}
\newacronym{LOHS}{LOHS}{Loss of Heat Sink}
\newacronym{LOLA}{LOLA}{Loss of Large Area}
\newacronym{LP}{LP}{linear program}
\newacronym{MA}{MA}{minor actinide}
\newacronym{MCNP}{MCNP}{Monte Carlo N-Particle code}
\newacronym{MILP}{MILP}{mixed-integer linear program}
\newacronym{MIT}{MIT}{the Massachusetts Institute of Technology}
\newacronym{MOAB}{MOAB}{Mesh-Oriented datABase}
\newacronym{MOOSE}{MOOSE}{Multiphysics Object-Oriented Simulation Environment}
\newacronym{MOX}{MOX}{mixed oxide}
\newacronym{MSBR}{MSBR}{Molten Salt Breeder Reactor}
\newacronym{MSRE}{MSRE}{Molten Salt Reactor Experiment}
\newacronym{MSR}{MSR}{Molten Salt Reactor}
\newacronym{NAGRA}{NAGRA}{National Cooperative for the Disposal of Radioactive Waste}
\newacronym{NEAMS}{NEAMS}{Nuclear Engineering Advanced Modeling and Simulation}
\newacronym{NEUP}{NEUP}{Nuclear Energy University Programs}
\newacronym{NFCSim}{NFCSim}{Nuclear Fuel Cycle Simulator}
\newacronym{NGNP}{NGNP}{Next Generation Nuclear Plant}
\newacronym{NMWPC}{NMWPC}{Nuclear MW Per Capita}
\newacronym{NNSA}{NNSA}{National Nuclear Security Administration}
\newacronym{NPRE}{NPRE}{Department of Nuclear, Plasma, and Radiological Engineering}
\newacronym{NQA1}{NQA-1}{Nuclear Quality Assurance - 1}
\newacronym{NRC}{NRC}{Nuclear Regulatory Commission}
\newacronym{NSF}{NSF}{National Science Foundation}
\newacronym{NSSC}{NSSC}{Nuclear Science and Security Consortium}
\newacronym{NUWASTE}{NUWASTE}{Nuclear Waste Assessment System for Technical Evaluation}
\newacronym{NWF}{NWF}{Nuclear Waste Fund}
\newacronym{NWTRB}{NWTRB}{Nuclear Waste Technical Review Board}
\newacronym{OCRWM}{OCRWM}{Office of Civilian Radioactive Waste Management}
\newacronym{ORION}{ORION}{ORION}
\newacronym{ORNL}{ORNL}{Oak Ridge National Laboratory}
\newacronym{PARCS}{PARCS}{Purdue Advanced Reactor Core Simulator}
\newacronym{PBAHTR}{PB-AHTR}{Pebble Bed Advanced High Temperature Reactor}
\newacronym{PBFHR}{PB-FHR}{Pebble-Bed Fluoride-Salt-Cooled High-Temperature Reactor}
\newacronym{PEI}{PEI}{Peak Environmental Impact}
\newacronym{PH}{PRONGHORN}{PRONGHORN}
\newacronym{PRKE}{PRKE}{Point Reactor Kinetics Equations}
\newacronym{PSPG}{PSPG}{Pressure-Stabilizing/Petrov-Galerkin}
\newacronym{PWAR}{PWAR}{Pratt and Whitney Aircraft Reactor}
\newacronym{PWR}{PWR}{Pressurized Water Reactor}
\newacronym{PyNE}{PyNE}{Python toolkit for Nuclear Engineering}
\newacronym{PyRK}{PyRK}{Python for Reactor Kinetics}
\newacronym{QA}{QA}{quality assurance}
\newacronym{RDD}{RD\&D}{Research Development and Demonstration}
\newacronym{RD}{R\&D}{Research and Development}
\newacronym{RELAP}{RELAP}{Reactor Excursion and Leak Analysis Program}
\newacronym{RIA}{RIA}{Reactivity Insertion Accident}
\newacronym{RIF}{RIF}{Region-Institution-Facility}
\newacronym{SFR}{SFR}{Sodium-Cooled Fast Reactor}
\newacronym{SINDAG}{SINDA{\textbackslash}G}{Systems Improved Numerical Differencing Analyzer $\backslash$ Gaski}
\newacronym{SKB}{SKB}{Svensk K\"{a}rnbr\"{a}nslehantering AB}
\newacronym{SNF}{SNF}{spent nuclear fuel}
\newacronym{SNL}{SNL}{Sandia National Laboratory}
\newacronym{STC}{STC}{specific temperature change}
\newacronym{SUPG}{SUPG}{Streamline-Upwind/Petrov-Galerkin}
\newacronym{SWF}{SWF}{Separations and Waste Forms}
\newacronym{SWU}{SWU}{Separative Work Unit}
\newacronym{TRIGA}{TRIGA}{Training Research Isotope General Atomic}
\newacronym{TRISO}{TRISO}{Tristructural Isotropic}
\newacronym{TSM}{TSM}{Total System Model}
\newacronym{TSPA}{TSPA}{Total System Performance Assessment for the Yucca Mountain License Application}
\newacronym{ThOX}{ThOX}{thorium oxide}
\newacronym{UFD}{UFD}{Used Fuel Disposition}
\newacronym{UML}{UML}{Unified Modeling Language}
\newacronym{UOX}{UOX}{uranium oxide}
\newacronym{UQ}{UQ}{uncertainty quantification}
\newacronym{US}{US}{United States}
\newacronym{UW}{UW}{University of Wisconsin}
\newacronym{VISION}{VISION}{the Verifiable Fuel Cycle Simulation Model}
\newacronym{VV}{V\&V}{verification and validation}
\newacronym{WIPP}{WIPP}{Waste Isolation Pilot Plant}
\newacronym{YMR}{YMR}{Yucca Mountain Repository Site}


\usepackage{algorithm}
\usepackage{algpseudocode}
\usepackage{array}

\usepackage{xspace}
\usepackage{multirow}

\newcommand{\cycamore}{\textsc{Cycamore}\xspace}
\newcommand{\cyclus}{\textsc{Cyclus}\xspace}
\newcommand{\SN}{S$_N$}
\renewcommand{\vec}[1]{\bm{#1}} %vector is bold italic
\newcommand{\vd}{\bm{\cdot}} % slightly bold vector dot
\newcommand{\grad}{\vec{\nabla}} % gradient
\newcommand{\ud}{\mathop{}\!\mathrm{d}} % upright derivative symbol

\begin{document}
%%%%%%%%%%%%%%%%%%%%%%%%%%%%%%%%%%%%%%%%%%%%%%%%%%%%%%%%%%%%%%%%%%%%%%%%%%%%%%%%
\section{Introduction}

In 2020, a \gls{haleu} workshop report led by Monica Regalbuto \cite{regalbuto_high_assay_2020} highlighted the unique regulatory challenges of establishing a \gls{haleu} fuel cycle in the \gls{usa}. It noted that part of enriching \gls{haleu} is first to produce \gls{leup}, defined as between 5\% and 10\% $^{235}U$ enrichment. The report notes that \gls{leup} facilities would fall under a similar category of regulations as our existing \gls{leu} fuel cycle, a significant advantage for the \gls{nrc} and the United States government. Since then, the idea has cropped up in several forums that, if a reactor could be redesigned to accommodate \gls{leup}, it could serve as a useful intermediary step to \gls{haleu}.

% To enrich to up to \gls{haleu}, facilities such as TRISO-X LLC and Kairos Power Atlas Fuel Fabrication Facility must build . Thus, \gls{leup} is an attractive intermediary step for servicers wishing to minimize the size of a Category II facility (thereby reducing costs) as it is the same category we have historically licensed for \gls{leu} fuel enrichment.

% This report also notes that \gls{leup} could be used in existing \gls{leu} reactors, which would allow for a more rapid deployment of \gls{haleu} in the \gls{us}. This is an important point, as it would allow for a more rapid transition to \gls{haleu} without the need for new reactor designs or modifications to existing reactors.

% One of the primary advantages for a fuel cycle containing \gls{leup} is that the facility to produce it would fall under the same licensing category as \gls{leu} fuel. The \gls{nrc} defines a \textit{special nuclear material of low strategic significance} as meeting one of three criteria, the most notable of which for our purposes is "(3) 10,000 grams or more of uranium-235 (contained in uranium enriched above natural but less than 10 percent in the U–235 isotope)," \cite{nrc_catiii}. This facility definition is where the upper limit of the \gls{leup} range arises.

In this work, we use \cyclus to model the \gls{nfc} of a fuel cycle that deploys Westinghouse AP1000s, \gls{xe}, and \gls{usnc} \gls{mmr}. After 10 years of operation with \gls{leup} fuel, the advanced reactor fleet transitions to deploying \gls{haleu} versions of the \gls{xe} and \gls{mmr}. We will evaluate the impact of this transition on the delayed \gls{swu} and mass of \gls{haleu} required to meet various demand growth scenarios.


%%%%%%%%%%%%%%%%%%%%%%%%%%%%%%%%%%%%%%%%%%%%%%%%%%%%%%%%%%%%%%%%%%%%%%%%%%%%%%%%
\section{\cyclus}

\cyclus \cite{huff_cyclus_intro_2016} is an agent-based \gls{nfc} simulator
that is versatile, open-source, and modular. The software achieves this
versatility through a series of generic archetypes that are primarily
transaction-based. Over the years, the user community and developers have
created a litany of nuclear-specific archetypes for everything from
proliferation assessment to fuel burnup. Many standard fuel cycle facilities
have been implemented in the \cycamore repository \cite{Carlsen_cycamore_2014},
which holds technology-agnostic archetypes for material sources, material
sinks, enrichment services, separations capabilities, storage services, and a
generic reactor.

% \begin{figure}[!htbp]
%   \centering
%   \includegraphics[width=0.4\textwidth]{../images/cyclus_logo.png}
%   \label{fig:cyclus_logo}
% \end{figure}

\section{Scenarios}
We will deploy reactors to meet energy demands drawn from \gls{eia} and \gls{doe} projections cited in Table \ref{tab:demand_scenarios}. The Low growth scenarios correspond to 5, 10, and 15\% increases in energy demand by 2050, and the High growth scenarios correspond to a doubling and a tripling of nuclear energy by 2050.

\begin{table}[h]
  \centering
  \caption{Demand Growth Scenarios}
  \label{tab:demand_scenarios}
  \begin{tabular}{l c c}
      \hline
      \textbf{Demand Growth} & \textbf{Year-to-Year Increase} & \textbf{Source}\\
      \hline
      No Growth & 0\% & na\\
      \multirow{3}{*}{Low Growth} & 0.17\% & \cite{eia_aeo_2023}\\
       & 0.5\% & \cite{eia_aeo_2023}\\
       & 1\% & \cite{eia_aeo_2023}\\
       \multirow{2}{*}{High Growth} & 3.5\% & \cite{julie_liftoff_pathways_2024} \\
       & 5.6\% & \cite{julie_liftoff_pathways_2024}\\
      \hline
  \end{tabular}
\end{table}

In Table \ref{tab:deployment_schemes}, we describe the three schemes we used to deploy reactors to meet the energy demand. Through the greedy scheme, we are not attempting to capture the complexity of the deployment problem but rather to explore the implications of deploying a certain number of reactors. This random deployment is a proxy for the complexity of the real-world deployment problem; however, it does not include the nuance of how individual deployments meet an end user's needs, which will drive the strategic decisions that utilities and ratepayers behind the meter make in their reactor choices. Combining the random and greedy schemes allows us to inject some uncertainty around which reactor will be deployed at any given time while ensuring meeting the demand efficiently.

\begin{table}[h]
  \centering
  \caption{Deployment Schemes}
  \label{tab:deployment_schemes}
  \begin{tabular}{>{\raggedright}p{0.24\linewidth}>{\raggedright\arraybackslash}p{0.60\linewidth}}
      \hline
      Scheme & Description \\
      \hline
      Greedy Deployment & Deploy the largest
      reactor first at each time step, fill in the remaining capacity with
      the next smallest, and so on. \\
      Random Deployment & Use a date and hour as seed to sample the
      reactors list randomly. \\
      Initially Random, Greedy Deployment & Randomly deploy reactors until
      a reactor bigger than the remaining capacity is proposed for each time step,
      then fills the remaining capacity with the greedy algorithm. \\
      \hline
  \end{tabular}
\end{table}

% \begin{algorithm}
% \caption{Greedy Reactor Deployment Algorithm}
% \begin{algorithmic}[1]
%     \State Initialize demand
%     \While{demand exists}
%         \State Select the largest reactor that does not exceed demand
%         \State Deploy reactors until the next reactor exceeds demand
%         \State Update demand
%     \EndWhile
% \end{algorithmic}
% \end{algorithm}

% \begin{algorithm}
%   \caption{Random Reactor Deployment Algorithm}
%   \begin{algorithmic}[1]
%       \State Initialize demand
%       \While{demand exists}
%           \State Randomly deploy a reactor that does not exceed demand
%           \State Update demand
%       \EndWhile
%   \end{algorithmic}
%   \end{algorithm}

%   \begin{algorithm}
%     \caption{Random + Greedy Reactor Deployment Algorithm}
%     \begin{algorithmic}[1]
%         \State Initialize demand
%         \While{demand exists}
%             \State Randomly deploy a reactor
%             \If{demand is exceeded}
%                 \State Remove last reactor
%                 \If{demand still exists}
%                     \State Select the largest reactor that does not exceed demand
%                     \State Deploy until the next reactor exceeds demand
%                     \State Update demand
%                 \EndIf
%             \EndIf
%         \EndWhile
%     \end{algorithmic}
%     \end{algorithm}

% \begin{subequations} \label{eqs:fullTransport}
% \begin{multline} \label{eq:fullTransportVol}
%   \vec{\Omega}\vd \grad \psi(\vec{x}, \vec{\Omega})
%   + \sigma(\vec{x}) \psi (\vec{x}, \vec{\Omega})
% \\ =
%   \frac{\sigma_s(\vec{x})}{4\pi} \int_{4\pi} \psi(\vec{x},\vec{\Omega}')
%   \ud\Omega' + \frac{q(\vec{x})}{4\pi}
%   \equiv \frac{1}{4\pi} Q(\vec{x}) \,,
% \end{multline}
% inside $\vec{x} \in V$, $\vec{\Omega} \in 4\pi$, with an incident boundary
% condition
% \begin{equation} \label{eq:fullTransportBndy}
%   \psi(\vec{x}, \vec{\Omega}) = \psi^b(\vec{x}, \vec{\Omega}) \,,
%  \quad \vec{x} \in \partial V, \ \vec{\Omega} \vd \vec{n} < 0\,.
% \end{equation}
% \end{subequations}

%%%%%%%%%%%%%%%%%%%%%%%%%%%%%%%%%%%%%%%%%%%%%%%%%%%%%%%%%%%%%%%%%%%%%%%%%%%%%%%%
% \section{Results and Analysis}
% Table \ref{tab:enrichment_levels} shows the various levels of enrichment for uranium that we will use in this work.

% \begin{table}[!htbp]
%    \centering
%    \caption{Enrichment levels and their ranges.}
%    \label{tab:enrichment_levels}
%    \begin{tabular}{c c}
%       \hline
%       \textbf{Enrichment Level} & \textbf{Range [\%  $^{235}$U]} \\
%       \hline
%       Natural & < 0.711 \\
%       \gls{leu} & 0.711-5 \\
%       \gls{leup} & 5-10 \\
%       \gls{haleu} & 10-20 \\
%       % \gls{heu} & $\geq$ 20  \\
%       \hline
%    \end{tabular}
% \end{table}

%%%%%%%%%%%%%%%%%%%%%%%%%%%%%%%%%%%%%%%%%%%%%%%%%%%%%%%%%%%%%%%%%%%%%%%%%%%%%%%%
% \subsection{Subsection Goes Here (Heading B)}
% The user must manually capitalize initial letters of a subsection heading.

% For those who like equations in their papers, \LaTeX\ is a good choice. Here is
% an equation for the Marshak diffusion boundary condition:
% \begin{equation} \label{eq:marshak}
%   4 J^- = \phi + 2 D \vec{n} \vd \grad \phi \,.
% \end{equation}
% If we so choose, we can effortlessly reference the equation later.

% Another paragraph starts with Eq.~\eqref{eq:marshak} and sets $J^-$ to zero, a
% vacuum boundary condition:
% \begin{equation*}
%   0 = \phi + \frac{2}{3} \frac{1}{\sigma} \vec{n} \vd \grad \phi \,.
% \end{equation*}
% The extrapolation distance is $2/3$. A more detailed asymptotic analysis yields
% an extrapolation distance of about $0.71045$.



% Later on, we can include a table, even one that spans two columns such as
% Table~\ref{tab:widetable}.
%%%%%%%%%%%%%%%%%%%%%%%%%%%%%%%%%%%%%%%%
% \begin{table*}[htb]
%   \centering
%   \caption{Example of a Really Wide Table that Might Not Normally Fit in the Document}
%   \begin{tabular}{llllllllll}\toprule
%       & $\phi_T(0)$      & $\phi_T(10)$      & $\phi_T(20)$      &
%       $\phi_D(0)$      & $\phi_D(10)$      & $\phi_D(20)$      & $\rho$      &
%       $\varepsilon$      & $N_\text{it}$
% \\ \midrule
% $c=0.999$  & 0.9038 & 20.63 & 31.24 & 0.9087 & 20.63 & 31.23 & 0.2192 & $10^{-7}$ & 15
% \\
% $c=0.990$  & 0.3675 & 13.04 & 24.7 & 0.3696 & 13.04 & 24.69 & 0.2184 & $10^{-7}$ & 15
% \\
% $c=0.900$  & 0.009909 & 4.776 & 17.64 & 0.009984 & 4.786 & 17.63 & 0.2118 & $10^{-7}$ & 14
% \\
% $c=0.500$  & $6.069\times 10^{-5}$ & 2.212 & 15.53 & 6.213$\times 10^{-5}$ & 2.239 & 15.53 & 0.2068 & $10^{-7}$ & 13
% \\
% \bottomrule
% \end{tabular}
%   \label{tab:widetable}
% \end{table*}
%%%%%%%%%%%%%%%%%%%%%%%%%%%%%%%%%%%%%%%%
% Notice how the table reference uses a Roman numeral
% for its numbering scheme, whereas the figure reference uses an Arabic numeral.
% For one-column tables, use the \verb|table| environment; two-column tables use
% \verb|table*|. The same applies to figures.

%%%%%%%%%%%%%%%%%%%%%%%%%%%%%%%%%%%%%%%%%%%%%%%%%%%%%%%%%%%%%%%%%%%%%%%%%%%%%%%%
% \section{Conclusions (Heading A)}

% The included ANS style file and this clear example file are a panacea for
% the hours of headache that invariably results from formatting a document in
% Microsoft Word.

%%%%%%%%%%%%%%%%%%%%%%%%%%%%%%%%%%%%%%%%%%%%%%%%%%%%%%%%%%%%%%%%%%%%%%%%%%%%%%%%
\appendix
% \section{Appendix}

% Numbering in the appendix is different:
% \begin{equation} \label{eq:appendix}
%   2 + 2 = 5\,.
% \end{equation}
% and another equation:
% \begin{equation} \label{eq:appendix2}
%   a + b = c\,.
% \end{equation}

%%%%%%%%%%%%%%%%%%%%%%%%%%%%%%%%%%%%%%%%%%%%%%%%%%%%%%%%%%%%%%%%%%%%%%%%%%%%%%%%
\section{Acknowledgments}
This research was performed, in part, using funding received from the DOE
Office of Nuclear Energy's Nuclear Energy University Program (Project 23-29656
DE-NE0009390) 'Illuminating Emerging Supply Chain and Waste Management
Challenges'.

%%%%%%%%%%%%%%%%%%%%%%%%%%%%%%%%%%%%%%%%%%%%%%%%%%%%%%%%%%%%%%%%%%%%%%%%%%%%%%%%
\bibliographystyle{ans}
\bibliography{../bibliography}
\end{document}

